\section*{ВИСНОВКИ}
\addcontentsline{toc}{section}{ВИСНОВКИ}

У ході виконання курсової роботи було розроблено програмну систему діалогової взаємодії з великими мовними моделями з web- та консольним інтерфейсами. Система забезпечує створення й керування чатами, надсилання повідомлень, отримання відповідей від локальної LLM-моделі, збереження історії у PostgreSQL та роботу з локальними моделями через Hugging Face Transformers.

У першому розділі було проаналізовано предметну область LLM-систем, розглянуто основні підходи до побудови діалогових застосунків і виконано порівняння з існуючими рішеннями, такими як ChatGPT, Claude, Google Gemini, LM Studio, Ollama та Open WebUI. На основі аналізу обґрунтовано доцільність створення власної навчальної модульної системи.

У другому розділі було спроєктовано архітектуру програмного продукту. Визначено основні рівні системи: web- та CLI-інтерфейси, бізнес-логіка, сховище даних, LLM-рівень і конфігураційний рівень. Також сформовано функціональні й нефункціональні вимоги, описано структуру бази даних і наведено UML-діаграми системи.

У третьому розділі було описано реалізацію програмної системи. Основну бізнес-логіку зосереджено у \texttt{ChatManager}, збереження даних реалізовано через \texttt{PostgresStorage}, а LLM-рівень --- через \texttt{TransformersLLMService} та \texttt{LLMRuntime}. У web-інтерфейсі реалізовано streaming generation, зупинку генерації, пошук повідомлень, експорт діалогів, відображення Markdown, блоків коду та LaTeX-формул через локальний KaTeX.

Особливу увагу приділено розширюваності системи. Реалізовано підтримку runtime-перемикання моделей, generation presets, локального registry моделей та PEFT LoRA/QLoRA-адаптерів. Завдяки цьому система не прив'язана до однієї конкретної моделі та може бути використана для подальших експериментів з різними локальними LLM.

Працездатність системи перевірено автоматизованими тестами. Тестами покрито доменну логіку, web endpoints, storage-рівень, runtime-керування моделлю, побудову prompt-ів, конфігурацію, міграції та допоміжні скрипти. Повний набір тестів успішно виконується і підтверджує коректність основних сценаріїв роботи.

У результаті поставлену мету досягнуто. Розроблена система є працездатним навчальним програмним продуктом, який демонструє повний цикл побудови локальної LLM-системи: від проєктування архітектури й бази даних до реалізації web-інтерфейсу, інтеграції локальної моделі, streaming-генерації та автоматизованого тестування.

Подальший розвиток системи може включати покращення якості відповідей шляхом підбору сильніших локальних моделей, додавання повноцінної довготривалої пам'яті між чатами, розширення механізму адаптерів, покращення UI для керування сховищем та впровадження авторизації у разі переходу від локального навчального використання до багатокористувацького сценарію.

\newpage
